\section{Proposed Solutions}\label{sec:solutions}
% Author wording (2026-11-27), kept maximally; only factual errors corrected.
% Pointers as in Sec. III. RC = contracts/test/termpaper/OrderBookMitigated.sol (research copy, NOT the engine).

We describe each issue using change, functions, status, effect and what it leaves open. Table~\ref{tab:mitigations} provides a summary.
Only M1 and M2 were implemented, in the research copy. The submitted engine contains neither. Implementing M3 to clear crossings before a new order takes is future work. % RC header 14-17; EM M3 (none)

\subsection{No Crossed Rest and a Minimum Order Size}\label{sec:m12}
\emph{Change.}
When a taker's fills run out, M2 returns the escrow of a remainder that is still crossable with the best order on the other side and closes it. % EM M2a; RC:75-78, 162-164
M1 rejects orders less than the \texttt{minOrder} parameter. % EM M1a; RC:52-53, 66, 98
\emph{Code impact.}
M2 adds one view function, \texttt{\_stillCrossed}, called in \texttt{placeLimitOrder} and in the activation step of \texttt{triggerStopOrders}. M1 adds one check to each placement function and reuses the existing error \texttt{ZeroAmount}. % RC:66, 75-78, 81-88, 98, 162-164
\emph{Status.} Both changes are implemented in the research copy of the engine only. % RC header

\emph{Measured effect.}
With M2, issue 007 no longer pays: Carol resubmits her refunded remainder and pays 14.99999999999 quote for 10 base, and Mallory ends $1.499\times10^{-8}$ quote up and $10^{-8}$ base down. This represents a net loss of $10^{-11}$ to Mallory when her base is valued at 1.5 (E5). % EV E5 table; TermPaper.t.sol:309; EM M2a
The stop snipe also no longer pays, and the book is left uncrossed (E7). % EV E7 table; EM M2b
In the one book shape we measured, M2 added 606, 26 and 1{,}071 gas to three calls respectively. % EV E6 table; EM M2c

\emph{What is still open.}
Neither of these changes on their own is sufficient. % EM M2a caveat, M1b
With M2 alone, Mallory can still repeat the ten dust asks before every resubmission. Each round gives Carol $10^{-8}$ base, so, extrapolating from the five rounds we ran, it would take $10^{9}$ rounds to fill 10 base, turning theft into indefinite delay (E10). % EV E10 row 1 (5 run, 1e9 extrapolated)
In the snipe, M2 cancels Carol's activated stop-loss after it sells $10^{-8}$ base and refunds the rest, leaving her with 9.99999999 base and no stop-loss in a falling market (E7). % EV E7 table and per-test review
M1 with a minimum of 0.1 base bounds the repeated attack, meaning Carol is filled in 10 rounds for 14.99 quote, and Mallory's position, valued at 1.5, loses 0.010 quote. This happens while Mallory spends 3.3 times the gas Carol spent (E10). % EV E10 row 2; CM griefing table; EM M1d
M1 does not reduce the stop hunt in E8, and it was not tested against the snipe. % EV E8; EV E7 (minOrder 0)
Only one minimum was tested, so the size required for a given token is unknown. A fixed minimum also excludes small traders, and E10 assumes that Mallory sees each resubmission. % EM M1a; EV E10 "Assumes"

\subsection{Hardened Stop Activation}\label{sec:m67}
\emph{Proposed change.}
Against the stop hunt, M6 would permit only fills above a minimum size to update the trigger reference, or would use a time-weighted average price to set the trigger. M7 would reject or clip a stop order if its limit is beyond a specified band from the trigger. % EM M6, M7
M6 changes \texttt{\_settleFill}, which writes \texttt{lastTradePrice}, and \texttt{\_eligibleStop}, which reads it. M7 changes \texttt{placeStopLimitOrder} to reject a stop, or the activation step of \texttt{triggerStopOrders} to clip it. % OB:242-248, 291, 73-85, 127-145
\emph{Status.} Neither of these changes is implemented or tested. % OPEN_QUESTIONS M6, M7
\emph{Known issues.}
E8 shows a size floor on orders does not resolve the issue, and we argue no floor would help, as a self-trade is essentially a wash trade that nets to zero. % EV E8; EM M1b; reasoning
As such, M6 should use a time-weighted average, or an average weighted by volume over a specified time window, so that no single trade sets the trigger.
Moosavi and Clark withhold a clearing price for a similar reason, noting that a miner could fix it with a single wash trade, and suggest a rolling median price as a countermeasure, which they borrow from Uniswap~\cite{src09moosavi2023}. % [09] p. 19 App. D (F10); EM C2g, M6
We believe M7 would cap Carol's loss at the band, but some stops could go unexecuted in a fast fall. % design reasoning (EM M7: none)

\subsection{Dust Wall, Tokens, Keepers and the Front End}\label{sec:m8to12}
For C4, M1 would increase the capital locked in a dust wall by the ratio of the minimum to one lot, which we did not test, and a price-level book (M8) would allow an insertion to skip over whole price levels. % EM M1, M8a
We therefore expect insertions to walk through price levels rather than single orders, touching \texttt{\_tryInsert} and changing the data structure. In Lissy, the data structure alone changed the gas cost to remove 50 entries about fivefold across five designs~\cite{src09moosavi2023}. % OB:320-343; [09] p. 4 to 5 Table 3 (F2); EM M8a
For C6, a keeper reward (M9) would need a trading fee, which the engine lacks, and we expect it to become a prize that trading bots bid for with fees, as in the priority gas auctions of Daian et al.~\cite{src01daian2020}. % [01] p. 1, p. 10 to 14 (F1, F4, analogy); EM M9
For C5, we expect pull payments in \texttt{\_settleFill}~\cite{src15grech2018} to stop a refusing token from blocking a fill. % [15] p. 6, p. 7 to 8 (F3, F5); EM M10
Alternatively, an allowlist in \texttt{MarketFactory} could restrict which tokens can be listed, like the protocol allowlists suggested by Zhou et al.~\cite{src06zhou2023} (M10); the cost is an extra withdrawal step or the loss of permissionless listing. % [06] p. 12 Insight 4 (F5); EM M10
For C1, the front end could use tighter default limits (M12), at the cost of orders that fill less often. % EM M12
E1 shows that a limit 0.1\% above the ask shrinks the loss to 0.015 quote, against 5.0 at 33\%, without removing it. % EV E1 rows
After Uniswap changed its slippage setting from a static default of 0.5\% to a dynamic range of 0.1\% to 5\% in March 2023, losses on that AMM per 100 USD of swapped value decreased by 54.7\%~\cite{src34chemaya2025}. % [34] p. 2, p. 10 (F3), p. 14 Table 5 (F5); EM M12

\subsection{Design Directions: Batching and Ordering Protection}\label{sec:m4511}
According to Budish et al., in a frequent batch auction (M5), all orders in an interval are cleared at the same price. Time priority then applies only across intervals. Therefore, a fast trader can snipe stale quotes in only a fraction $\delta/\tau$ of the time, with a speed advantage $\delta$ and an interval $\tau$~\cite{src11budish2015}. % [11] Sec. 7.1, 7.2 Prop. 8 (F6, F7), no pages; EM M5a
We argue that, with respect to C1 and C3, a batch auction would remove the need to engage in a time race to be first, and Moosavi and Clark argue that an on-chain call market resists basic insertion and displacement attacks~\cite{src09moosavi2023}. % [09] p. 8 Table 7 (F7); EM M5b
However, in M5, \texttt{placeLimitOrder} matching is replaced with a separate clearing call, resulting in the worst-case throughput described in Section~\ref{sec:motivation}. % [09] p. 6 to 7 Table 5 (F3); EM M5e
We found no production measurements to validate that batch DEXs reduce extractable value. % OPEN_QUESTIONS M5d (gap d)
Commit and reveal (M4) would add a commit and a reveal step prior to each \texttt{placeLimitOrder} and would incur the costs and limits described in Section~\ref{sec:related}. % EM M4a to M4c
As discussed in Section~\ref{sec:related}, chain-level protection (M11) needs no contract change but has similar trade-offs. As an example, the batched threshold encryption proposed by Choudhuri et al.\ took 2{,}818.6 ms per committee member, single-threaded, to compute a partial decryption of a 512-transaction batch~\cite{src29choudhuri2024}. % [29] p. 12 Table 2 (F5); EM M11d
None of these is implemented, and we argue that neither commit and reveal nor chain-level protection, used individually, prevents the one-transaction snipe, since it reorders no victim transaction: Carol's stop order is already public. % EV E9 atomic; IOrderBook.sol:49; reasoning

% Table II. Built from EVIDENCE_MAP.md section B. Strength: S strong, M moderate, W weak, none = no evidence.
% "E" = measured in our experiments; "L" = literature only. P2: check every cell against EVIDENCE_MAP before submission.
\begin{table}[t]
\caption{Mitigations against considerations C1 to C6, with evidence}
\label{tab:mitigations}
\centering
\footnotesize
\begin{tabular}{@{}lllll@{}}
\toprule
ID & Mitigation & Targets & Evidence & Strength \
\midrule
M1 & Minimum order size & C1, C3, C4 (not C2) & E5, E8, E10 & S \
M2 & No crossed rest & C1, C3 & E5, E6, E7, E10 & S \
M3 & Clear crossings before taking & C1, C3 & none & W \
M4 & Commit and reveal & C1 to C3 & L & M \
M5 & Frequent batch auction & C1 to C3 & L & M \
M6 & Hardened stop reference price & C2 & L (indirect) & W \
M7 & Trigger-to-limit band & C2, C3 & none & none \
M8 & Price-level book & C4 & L & M \
M9 & Keeper reward & C6 & L & M \
M10 & Token allowlist or pull payments & C5 & L & M \
M11 & Chain-level ordering protection & C1 to C3 & L & M to S \
M12 & Frontend guards & C1 & E1 (slack sweep) & W \
\bottomrule
\end{tabular}
\end{table}


\emph{Recommended order of adoption.}
We would implement M2 together with M1 first, since they are the only changes with measured impact. Changes M1 and M2 affect only the placement and stop-activation paths, and E10 shows that M2 without M1 turns theft into delay. % EV E5 to E7, E10; RC:66, 75-78, 98, 162-164
Next, we would implement M6 and M7 for stops, as M1 and M2 have no effect on C2, and the canceled stop of C3 needs a better activation rule; both need experiments we have not run. % EV E8, E7; OPEN_QUESTIONS M6, M7
Finally, we would implement M12 and M10, which we expect to reduce potential exposure without changing the matching rule. % EM M12, M10
Batching and ordering protection are larger redesigns, of the engine or of the chain, whose cost we can cite but have not measured. % EM M5e, M4c, M11d
